\documentclass[../../main.tex]{subfiles}

\begin{document}

\chapter[Teori Bahasa Formal dan Hierarki Chomsky]{Teori Bahasa Formal\\ dan Hierarki Chomsky}
\label{chap:bahasa-formal}

\begin{outcomebox}
\textbf{Sub-CPMK:} \subcpmk{1.4}\\[3pt]
Setelah menyelesaikan bab ini, mahasiswa mampu:
\begin{enumerate}[label=\arabic*.,leftmargin=*]
  \item mendefinisikan alfabet, string, bahasa, dan tata bahasa secara
        formal;
  \item membedakan tata bahasa regular, bebas konteks, sensitif
        konteks, dan tanpa batas;
  \item menjelaskan hubungan setiap kelas bahasa dengan mesin pengenal
        yang bersesuaian;
  \item membuktikan bahwa suatu bahasa tidak regular dengan lemma
        pompa;
  \item menjelaskan implikasi hierarki Chomsky bagi desain lexer,
        parser, dan analisis semantik.
\end{enumerate}
\end{outcomebox}

\section*{Materi Pokok}
\begin{itemize}[leftmargin=*]
  \item Definisi formal alfabet, string, bahasa, dan operasi bahasa.
  \item Tata bahasa, derivasi, bentuk kalimat, dan ekivalensi.
  \item Hierarki Chomsky Tipe 3 sampai Tipe 0.
  \item Otomata hingga, otomata pushdown, dan mesin Turing.
  \item Lemma pompa dan implikasi bagi front-end kompilator.
\end{itemize}

\subfile{section-01}
\subfile{section-02}
\subfile{section-03}
\subfile{section-04}
\subfile{section-05}

\section*{Contoh}
\addcontentsline{toc}{section}{Contoh}

\begin{contoh}[Mengklasifikasikan lima tata bahasa]
Tentukan tipe setiap tata bahasa berikut.
\begin{enumerate}[label=(\alph*),leftmargin=*]
  \item \(S \to aS \mid bS \mid \varepsilon\).
  \item \(S \to aSb \mid \varepsilon\).
  \item \(S \to aSBC \mid aBC\), \(CB \to BC\), \(aB \to ab\),
        \(bB \to bb\), \(bC \to bc\), \(cC \to cc\).
  \item \(S \to AB\), \(AB \to BA\), \(A \to a\), \(B \to b\).
  \item \(S \to Sa \mid b\).
\end{enumerate}
\textbf{Jawaban.}
\begin{enumerate}[label=(\alph*),leftmargin=*]
  \item Tipe 3 (regular), karena semua aturan berbentuk \(A \to aB\)
        atau \(A \to \varepsilon\).
  \item Tipe 2 (bebas konteks), karena aturan berbentuk
        \(A \to \alpha\) dan menghasilkan \(\{a^n b^n\}\).
  \item Tipe 1 (sensitif konteks), karena aturan seperti
        \(CB \to BC\) bergantung pada konteks dan panjang string tidak
        berkurang. Tata bahasa ini membangkitkan
        \(\{a^n b^n c^n\}\).
  \item Tipe 0 (tanpa batas), karena aturan \(AB \to BA\) tidak
        memenuhi batasan sensitif konteks.
  \item Tipe 3 (regular), karena aturan berbentuk \(A \to Ba\) atau
        \(A \to a\), yaitu left-linear.
\end{enumerate}
\end{contoh}

\begin{contoh}[Tabel pemetaan Tipe 3--0]
Tabel berikut merangkum pemetaan yang harus diisi pada Sub-CPMK 1.4.
\begin{center}
\small
\begin{tabular}{@{}c l l@{}}
\toprule
\textbf{Tipe} & \textbf{Automata} & \textbf{Contoh bahasa} \\
\midrule
3 & DFA/NFA & \(a^{*}b^{*}\) \\
2 & PDA & \(\{a^n b^n \mid n \ge 0\}\) \\
1 & Otomata terbatas linear & \(\{a^n b^n c^n \mid n \ge 1\}\) \\
0 & Mesin Turing & bahasa yang dikenali mesin Turing \\
\bottomrule
\end{tabular}
\end{center}
Setiap baris harus konsisten: contoh bahasa pada kolom terakhir harus
benar-benar dikenali oleh automata pada kolom kedua.
\end{contoh}

\begin{contoh}[Dari ekspresi reguler ke otomata]
Ekspresi reguler \((a|b)^{*}abb\) menyatakan bahasa seluruh string
atas \(\{a,b\}\) yang berakhir dengan \(abb\). Otomata hingga untuk
bahasa ini memiliki empat state yang mengingat prefiks terpanjang
yang cocok: \(\varepsilon\), \(a\), \(ab\), dan \(abb\). State
terakhir adalah state penerima. Setiap simbol baru memperbarui state
sesuai prefiks terpanjang, sehingga otomata hanya memerlukan memori
konstan.
\end{contoh}

\section*{Aktivitas Pembelajaran dan Praktikum}
\addcontentsline{toc}{section}{Aktivitas Pembelajaran dan Praktikum}

\begin{table}[H]
\centering
\caption{Rencana aktivitas Bab IV}
\label{tab:ch04-aktivitas}
\begin{tabularx}{\textwidth}{@{}llX@{}}
\toprule
\textbf{Sesi} & \textbf{Bentuk} & \textbf{Kegiatan} \\
\midrule
1 & Ceramah & Definisi formal bahasa dan tata bahasa \\
2 & Latihan & Klasifikasi tata bahasa pada hierarki Chomsky \\
3 & Diskusi & Hubungan tata bahasa, automata, dan fase kompilasi \\
4 & Lokakarya & Membuktikan bahasa tidak regular dan merancang PDA \\
\bottomrule
\end{tabularx}
\end{table}

\subsection*{Lokakarya Klasifikasi}
\begin{enumerate}[label=\arabic*.,leftmargin=*]
  \item Diberikan lima tata bahasa, tentukan tipenya dan jelaskan
        alasannya.
  \item Pilih satu bahasa regular dan satu bahasa bebas konteks.
        Tuliskan tata bahasa, automata, dan contoh string untuk
        masing-masing.
  \item Buktikan dengan lemma pompa bahwa
        \(\{a^n b^n \mid n \ge 0\}\) bukan regular.
  \item Rancang PDA untuk \(\{a^n b^{2n} \mid n \ge 0\}\).
  \item Diskusikan mengapa analisis semantik tidak cukup dilakukan
        dengan CFG.
\end{enumerate}

Bab ini tidak menambahkan increment proyek. Materi digunakan sebagai
dasar teoretis untuk Bab V--VII (leksikal) dan Bab VIII--XIII
(sintaksis).

\section*{Latihan dan Refleksi}
\addcontentsline{toc}{section}{Latihan dan Refleksi}

\subsection*{Latihan}
\begin{enumerate}[label=\arabic*.,leftmargin=*]
  \item Definisikan tata bahasa regular dan berikan dua contoh.
  \item Jelaskan perbedaan tata bahasa bebas konteks dan sensitif
        konteks.
  \item Tuliskan tata bahasa untuk \(\{a^n b^{2n} \mid n \ge 0\}\).
  \item Lengkapi tabel pemetaan Tipe 3--0 ke jenis automata beserta
        satu contoh bahasa untuk setiap tipe.
  \item Diberikan lima grammar, klasifikasikan setiap grammar dan
        jelaskan alasan klasifikasinya.
  \item Jelaskan mengapa mesin Turing lebih kuat daripada PDA.
  \item Sebutkan tiga sifat tertutup bahasa regular dan dua sifat
        yang tidak dimiliki bahasa bebas konteks.
  \item Mengapa lexer tidak memerlukan stack, sedangkan parser
        memerlukannya?
\end{enumerate}

\subsection*{Refleksi}
\begin{reflectionbox}
\begin{enumerate}[label=\arabic*.,leftmargin=*]
  \item Konsep hierarki Chomsky mana yang paling jelas bagi Anda?
  \item Apakah Anda dapat menjelaskan hubungan setiap tipe grammar
        dengan jenis automata yang bersesuaian?
  \item Bagaimana pemahaman hierarki ini mengubah cara Anda memandang
        token dan pohon sintaksis?
\end{enumerate}
\end{reflectionbox}

\section*{Asesmen}
\addcontentsline{toc}{section}{Asesmen}

Asesmen Bab IV mencakup Kuis Konsep (3\%) dan UTS Praktik (14\%).
Fokusnya adalah ketepatan tabel pemetaan Tipe 3--0, klasifikasi lima
grammar, dan hubungan teori dengan fase kompilasi.

\begin{table}[H]
\centering
\caption{Instrumen asesmen Bab IV}
\label{tab:ch04-asesmen}
\begin{tabularx}{\textwidth}{@{}lXl@{}}
\toprule
\textbf{Instrumen} & \textbf{Indikator} & \textbf{Bobot} \\
\midrule
Kuis Konsep & Definisi formal, hierarki Chomsky, dan jenis automata &
3\% (bagian) \\
UTS Praktik & Tabel pemetaan Tipe 3--0 dan klasifikasi lima grammar
beserta alasannya & 14\% (bagian) \\
Lokakarya & Tabel, klasifikasi, dan hasil uji lima grammar &
prasyarat \\
\bottomrule
\end{tabularx}
\end{table}

Kriteria kelulusan: mahasiswa dapat mengisi tabel pemetaan Tipe 3--0
ke automata, mengklasifikasikan lima grammar dengan alasan yang tepat,
dan menjelaskan mengapa setiap fase kompilator memerlukan tingkat
formalisme yang berbeda.

\section*{Checklist Pencapaian Kompetensi}
\addcontentsline{toc}{section}{Checklist Pencapaian Kompetensi}

\begin{checklist}
  \item Saya dapat mendefinisikan alfabet, string, dan bahasa.
  \item Saya dapat menulis tata bahasa dalam bentuk empat-tupel.
  \item Saya dapat menghasilkan derivasi leftmost dan rightmost.
  \item Saya dapat membedakan tipe 3, 2, 1, dan 0.
  \item Saya dapat menghubungkan setiap tipe dengan mesin pengenal.
  \item Saya dapat menggunakan lemma pompa untuk membuktikan bahasa
        tidak regular.
  \item Saya dapat merancang PDA sederhana.
  \item Saya dapat menjelaskan implikasi hierarki bagi lexer, parser,
        dan analisis semantik.
\end{checklist}

\section*{Rangkuman}
\addcontentsline{toc}{section}{Rangkuman}

Pada bab ini telah dibahas:
\begin{itemize}[leftmargin=*]
  \item definisi formal alfabet, string, bahasa, dan tata bahasa;
  \item derivasi, bentuk kalimat, ekivalensi, dan bentuk normal;
  \item hierarki Chomsky Tipe 3 sampai Tipe 0;
  \item otomata hingga, otomata pushdown, otomata terbatas linear,
        dan mesin Turing;
  \item tabel pemetaan Tipe 3--0 dan uji klasifikasi lima grammar;
  \item implikasi hierarki bagi front-end kompilator.
\end{itemize}

Bab V memulai perjalanan implementasi dengan analisis leksikal. Kita
akan melihat bagaimana ekspresi reguler dipakai untuk mendefinisikan
token, dan bagaimana teori Tipe 3 berubah menjadi \eng{lexer} yang
bekerja.

\section*{Evaluasi}
\addcontentsline{toc}{section}{Evaluasi}

\begin{enumerate}[label=\arabic*.,leftmargin=*]
  \item Tuliskan definisi formal tata bahasa dan jelaskan setiap
        komponennya.
  \item Klasifikasikan tata bahasa \(S \to aSb \mid \varepsilon\) dan
        jelaskan alasannya.
  \item Buktikan bahwa \(\{a^n b^n \mid n \ge 0\}\) bukan bahasa
        regular.
  \item Jelaskan mengapa \(\{a^n b^n c^n \mid n \ge 1\}\) bukan bahasa
        bebas konteks.
  \item Gambarkan PDA untuk \(\{a^n b^{2n} \mid n \ge 0\}\).
  \item Bandingkan daya ekspresif DFA, PDA, dan mesin Turing.
  \item Mengapa lexer cukup memakai otomata hingga, sedangkan parser
        memerlukan otomata pushdown?
  \item Berikan satu contoh kesalahan yang tidak dapat dideteksi oleh
        tata bahasa bebas konteks.
\end{enumerate}

\subsection*{Kunci dan Rubrik Singkat}
\begin{itemize}[leftmargin=*]
  \item Soal 1--2: ketepatan definisi dan klasifikasi (70\%),
        penjelasan (30\%).
  \item Soal 3--4: struktur pembuktian (60\%), ketepatan argumen
        (40\%).
  \item Soal 5: kebenaran transisi dan stack (80\%), kelengkapan
        (20\%).
  \item Soal 6--8: ketepatan konsep (60\%), contoh (40\%).
\end{itemize}


\ifSubfilesClassLoaded{%
  \bibliographystyle{apalike}%
  \bibliography{../../references}%
}{}

\end{document}
